
\section{分次模对应的拟凝聚层}

我们现在来描述\trProjective{} $A$-\trScheme{}上的\trQCSheaf。回想：\trProjective{} $A$-\trScheme{}是由一个以 $\mathbb{Z}^{\ge 0}$ \trGrading{}的\trRing{} $S_\bullet$ 的数据产出的，其中 $S_0=A$，且 $S_+$ 是一个\trFinitelyGenerated{}的\trIdeal（「$A$ 上有限生成的\trGradedRing」，\S4.5.6）。所得的\trScheme{}记作 $\Proj S_\bullet$。

设 $M_\bullet$ 是一个分次 $S_\bullet$-\trModule，\emph{以 $\mathbb{Z}$ \trGrading}。（在阅读下一节时，你或许会奇怪我们为什么不以 $\mathbb{Z}^{\ge 0}$ 分次。你将会看到，这无关紧要：当我们由 $\Proj S_\bullet$ 上的一个\trQCSheaf{}造出一个 $M_\bullet$ 时，以 $\mathbb{Z}$ \trGrading{}会让事情更干净。）我们按如下方式定义\trQCSheaf{} $\widetilde{M_\bullet}$。（我会避免把它称作 $\widetilde{M}$，因为那可能与\trAffine{}情形相混淆；但 $\widetilde{M_\bullet}$ \emph{并}不以任何方式分次。）对每个\emph{正\trDegree{}}的\trHomogeneous{}元素 $f$，我们在\trDistinguishedOpen{} $\mathrm{D}(\mathrm{f})=\{p\ :\ \mathrm{f}(p)\ne 0\}$ 上定义一个\trQCSheaf{} $\widetilde{M_\bullet}(\mathrm{f})$ 为 $\widetilde{((M_\bullet)_f)_0}$——注意 $((M_\bullet)_f)_0$ 是一个 $((S_\bullet)_f)_0$-\trModule，并回想 $\mathrm{D}(\mathrm{f})$ 已被等同于 $\Spec((S_\bullet)_f)_0$（\cref{ex:d-plus-is-spec}）。正如在 \eqref{eq:spec-localized-zero-graded} 中那样，下标 $0$ 意指「次数 $0$ 分片」。

$\widetilde{M_\bullet}(\mathrm{f})$ 与 $\widetilde{M_\bullet}(\mathrm{g})$ 到 $\mathrm{D}(\mathrm{fg})$ 上的限制之间存在显而易见的同构，满足\trCocycleCondition{}。（请你亲自把这件事想清楚，好确信你也同意「显而易见」这个词！）据\cref{ex:gluing-sheaves-cocycle}，这些层粘合为 $\Proj S_\bullet$ 上的单独一个层 $\widetilde{M_\bullet}$。此后我们便弃用临时记号 $\widetilde{M_\bullet}(\mathrm{f})$。

$\OO$-\trModule{} $\widetilde{M_\bullet}$ 显然是\trQC{}的，因为它在每个 $\mathrm{D}(\mathrm{f})$ 上都\trQC，而\trQC{}性是局部的。

\textbf{14.6.A.（练习）}\quad 给出 $\widetilde{M_\bullet}$ 在对应于\trHomogeneous\trPrimeIdeal{} $\mathfrak{p}\subset S_\bullet$ 的那个点处的\trStalk{}与 $((M_\bullet)_\mathfrak{p})_0$ 之间的一个\trIsomorphism。这里 $((M_\bullet)_\mathfrak{p})_0$ 意指「$(M_\bullet)_\mathfrak{p}$ 中\trDegree{}为 $0$ 的\trHomogeneous{}分式」，而 $\mathfrak{p}$ 应当是一个「相关」的\trHomogeneous\trPrimeIdeal（也就是说，$[\mathfrak{p}]$ 是 $\Proj S_\bullet$ 的一个点）。（注记：你可以利用这道练习，借「\trCompatible\trStalk」给出 $\widetilde{M_\bullet}$ 的另一个定义；参看\cref{ex:proj-sheaf-compatible-germs}。）

给定\trGradedModule{}之间的一个映射 $\phi\colon M_\bullet\to N_\bullet$，我们便得到层的一个诱导映射 $\widetilde{M_\bullet}\to\widetilde{N_\bullet}$。显式地说，在 $\mathrm{D}(\mathrm{f})$ 上，映射 $M_\bullet\to N_\bullet$ 诱导 $(M_\bullet)_f\to(N_\bullet)_f$，而后者又诱导 $\phi_f\colon((M_\bullet)_f)_0\to((N_\bullet)_f)_0$；并且这一构造在限制到更小的\trDistinguishedOpen{}时表现良好，也就是说，下图交换：
\[
\begin{tikzcd}[row sep=2.8em, column sep=3em]
((M_\bullet)_f)_0 \ar[r, "\phi_f"] \ar[d] & ((N_\bullet)_f)_0 \ar[d] \\
((M_\bullet)_{fg})_0 \ar[r, "\phi_{fg}"] & {((N_\bullet)_{fg})_0.}
\end{tikzcd}
\]

于是 $\sim$ 是一个从分次 $S_\bullet$-\trModule{}的\trCategory{}到 $\Proj S_\bullet$ 上\trQCSheaf{}的\trCategory{}的\trFunctor。

\textbf{14.6.B.（容易的练习）}\quad 证明：$\sim$ 是一个\trExact{}函子。（提示：这个构造里的一切都是正合的。）

\textbf{14.6.C.（练习）}\quad 证明：若 $M_\bullet$ 与 $M'_\bullet$ 在足够高的次数上一致，则 $\widetilde{M_\bullet}\cong\widetilde{M'_\bullet}$。然后证明：从分次 $S_\bullet$-\trModule（在相差\trIsomorphism{}的意义下）到 $\Proj S_\bullet$ 上的\trQCSheaf（在相差\trIsomorphism{}的意义下）的这个映射不是一个\trBijection。（说真的：请证明这并不是\trEquivalenceOfCategories。）

\textbf{14.6.D.（练习）}\quad 描述一个 $S_0$-\trModule{}的映射 $M_0\to\Gamma(\Proj S_\bullet,\widetilde{M_\bullet})$。（这预示了 \S15.7.5 中那个把分次模送到其\trSaturation{}的「\trSaturation{}映射」；见\mbox{练习~15.7.C}。）

\textbf{14.6.1.}\quad \mbox{练习~14.6.C} 表明 $\sim$ 并不是范畴的同构（或等价），但它已相去不远。这层关系类似于\trPresheaf{}与\trSheaf{}之间的关系，也类似于\trSheafification{}函子（见 \S15.7）。理解这件事的一个便利办法，是从\trFinitelyGenerated{}的 $S_\bullet$-\trModule{}之\trCategory{} $\mathit{f.g.Mod}_{S_\bullet}$ 出发，造出一个\trObject{}相同、但\trMorphism{}\emph{不同}的新\trCategory{} $\scr{C}$：在这个新\trCategory{}中，从 $M_\bullet$ 到 $N_\bullet$ 的\trMorphism{}是
\[
\Hom_{\scr{C}}(M_\bullet,N_\bullet):=\lim_{m\to\infty}\Bigl(\Hom_{\mathit{f.g.Mod}_{S_\bullet}}(M_{\bullet\ge m},N_{\bullet\ge m})\Bigr).
\]
（这里 $M_{\bullet\ge m}=\oplus_{n\ge m}M_n$\footnote{原书此处作「$\oplus_{n\ge m}M_\bullet$」，被加项当系 $M_n$（直和 $M_\bullet$ 自身不成义），径改。——译者注}，一如你所预期。）然后你可以证明：$\scr{C}$ 等价于 $\Proj S_\bullet$ 上\trFiniteType{}的\trQCSheaf{}的\trCategory，而这个等价对应于映射 $M_\bullet\mapsto\widetilde{M_\bullet}$。欢迎你去证明这些事实，但我们不会用到它们。

\textbf{14.6.2.\ 例子：$S_\bullet$ 的\trHomogeneous\trIdeal{}给出 $\Proj S_\bullet$ 的\trClosedSubscheme。}\quad 回想：\trHomogeneous\trIdeal{} $\mathrm{I}_\bullet\subset S_\bullet$ 给出一个\trClosedSubscheme{} $\Proj S_\bullet/\mathrm{I}_\bullet\hookrightarrow\Proj S_\bullet$。例如，设 $S_\bullet=k[w,x,y,z]$，于是 $\Proj S_\bullet\cong\mathbb{P}^3$。\trIdeal{} $\mathrm{I}_\bullet=(wz-xy,\ x^2-wy,\ y^2-xz)$ 给出的正是我们的老朋友——\trTwistedCubic（定义于\mbox{练习~9.3.A}）。

\textbf{14.6.E.（练习）}\quad 证明：若把函子 $\sim$ 施于分次 $S_\bullet$-\trModule{}的\trExactSequence
\[
0\longrightarrow\mathrm{I}_\bullet\longrightarrow S_\bullet\longrightarrow S_\bullet/\mathrm{I}_\bullet\longrightarrow 0
\]
我们便得到 $\Proj S_\bullet/\mathrm{I}_\bullet\hookrightarrow\Proj S_\bullet$ 的\trClosedSubscheme{}正合列 \eqref{eq:closed-subscheme-exact-sequence}。

$\Proj S_\bullet$ 的所有\trClosedSubscheme{}都以这种方式产生（见\mbox{练习~15.7.I}）。

\textbf{14.6.3.\ \trhRemark{}。}\quad 若 $M_\bullet$ 是\trFinitelyGenerated{}的（相应地，\trFinitelyPresented{}的、\trCoh{}的），则 $\widetilde{M_\bullet}$ 亦然 [Roh, (4.4.11)]。我们将用不到这个事实。（前两条的一个快速证明：若 $M_\bullet$ 是\trFinitelyGenerated{}的，相应地，\trFinitelyPresented{}的，则\trAffine\trCone{} $\Spec M_\bullet$ 上的\trQCSheaf{} $\widetilde{M_\bullet}$ 亦然。这两个概念都被\trPullback{}所保持，而 $\Proj S_\bullet$ 的诸\trStandardOpenSet{}是\trAffine\trCone{} $\Spec S_\bullet$ 的\trClosedSubscheme。）

